\documentclass[../DoAn.tex]{subfiles}
\begin{document}

\section{Yêu cầu và kiến trúc tổng quan}
\label{section:4.1}

\subsection{Sơ đồ use case}
\label{subsection:4.1.1}
\begin{figure}[H]
    \centering
    \includegraphics[width=0.7\textwidth]{Hinhve/usecase.png}
    \caption{Sơ đồ usecase tổng quan}
    \label{fig:usecase}
\end{figure}

\paragraph{Tác nhân tham gia}
Hệ thống nhận dạng biển số xe Việt Nam có một tác nhân duy nhất là Người dùng (User), đại diện cho bất kỳ cá nhân nào truy cập hệ thống qua giao diện web. Tác nhân này đóng hai vai trò khác nhau tùy theo trạng thái xác thực: với các chức năng quản lý tài khoản như đăng ký và đăng nhập, người dùng là khách chưa được xác thực; với các chức năng nghiệp vụ chính như xử lý video, giám sát và xem lịch sử, người dùng phải là thành viên đã đăng nhập hợp lệ. Hệ thống không phân biệt vai trò theo nhóm quyền (admin, operator...) mà kiểm soát truy cập hoàn toàn dựa trên phiên đăng nhập.

\paragraph{Các use case chính}
Hệ thống có bảy use case được mô tả trong sơ đồ, chia thành hai nhóm chức năng.

Nhóm quản lý tài khoản gồm ba use case. Đăng ký cho phép người dùng tạo tài khoản mới bằng thông tin cá nhân; hệ thống xác thực dữ liệu đầu vào, băm mật khẩu và lưu thông tin vào cơ sở dữ liệu. Đăng nhập cho phép người dùng có tài khoản xác thực danh tính và nhận JWT qua cookie bảo mật để truy cập các chức năng được bảo vệ. Đăng xuất hủy phiên làm việc hiện tại, xóa cookie và đưa người dùng về trạng thái khách.

Nhóm nghiệp vụ nhận dạng gồm bốn use case, chỉ khả dụng sau khi đăng nhập. Xử lý video nhận dạng biển số là use case trung tâm của hệ thống: người dùng tải video lên, backend khởi động pipeline ALPR đa tầng — phát hiện phương tiện, phát hiện biển số, tracking, đánh giá chất lượng crop, OCR và tổng hợp kết quả đa khung hình — rồi trả kết quả về frontend theo thời gian thực qua SSE. Giám sát và đánh dấu sự cố cho phép người dùng theo dõi luồng video liên tục (từ file upload hoặc nguồn RTSP) và đánh dấu một khoảng thời gian cụ thể để phân tích chuyên sâu, phục vụ các tình huống cần trích xuất bằng chứng. Xem kết quả và evidence nhận dạng cho phép người dùng xem lại kết quả của từng phiên xử lý, bao gồm danh sách biển số nhận dạng được, các xe bị từ chối (rejected vehicles) kèm lý do, và ảnh crop làm bằng chứng. Xem lịch sử nhận dạng cung cấp danh sách toàn bộ các phiên xử lý đã thực hiện, được lọc theo tài khoản đăng nhập, giúp người dùng tra cứu lại kết quả của các video đã xử lý trước đó.

\subsection{Thiết kế kiến trúc}
\paragraph{Lựa chọn kiến trúc phần mềm}
Hệ thống áp dụng kiến trúc ba lớp (Three-tier Architecture) kết hợp với mô hình Client–Server. Đây là kiến trúc phân tách ứng dụng thành ba tầng độc lập: tầng trình bày (Presentation Tier) chịu trách nhiệm hiển thị giao diện và tương tác với người dùng, tầng xử lý nghiệp vụ (Business Logic Tier) thực thi toàn bộ logic ứng dụng, và tầng dữ liệu (Data Tier) quản lý việc lưu trữ và truy vấn dữ liệu bền vững. Mỗi tầng chỉ giao tiếp với tầng liền kề thông qua giao diện được định nghĩa rõ ràng, giúp các tầng có thể thay đổi hoặc mở rộng độc lập mà không ảnh hưởng lẫn nhau.

\paragraph{Thiết kế tổng quan}
\begin{figure}[H]
    \centering
    \includegraphics[width=1.2\textwidth]{Hinhve/package.drawio.pdf}
    \caption{Biểu đồ phụ thuộc gói}
    \label{fig:Fig1}
\end{figure}

\paragraph{Tầng trình bày (Presentation Tier)}
Gồm ba package: Package Authentication chứa các màn hình xác thực người dùng, bao gồm Login (đăng nhập) và RouteGuard (bảo vệ định tuyến, ngăn người dùng chưa xác thực truy cập các trang nghiệp vụ). Package MainView tổ chức toàn bộ giao diện sau đăng nhập, gồm Dashboard (trang tổng quan), VideoUpload (tải video lên để xử lý), ResultView (xem kết quả và ảnh minh chứng của từng phiên nhận dạng), HistoryView (tra cứu lịch sử các phiên đã xử lý) và MonitorView (giám sát luồng video và đánh dấu sự kiện). Package APIClient đóng vai trò cầu nối duy nhất giữa tầng trình bày và tầng nghiệp vụ, gồm HTTPClient xử lý các yêu cầu REST đồng bộ và SSEClient nhận kết quả nhận dạng theo luồng thời gian thực từ server.
\paragraph{Tầng nghiệp vụ (Business Logic Tier)}
Gồm hai package. Package WebAPI tiếp nhận và xử lý toàn bộ yêu cầu HTTP từ client, được tổ chức thành bốn router theo chức năng: AuthAPI quản lý đăng ký, đăng nhập và xác thực JWT; RecognitionAPI tiếp nhận video đầu vào, khởi động pipeline xử lý và phát kết quả về client qua SSE; MonitorAPI quản lý luồng giám sát và các sự kiện được đánh dấu; HistoryAPI cung cấp dữ liệu lịch sử phiên nhận dạng theo tài khoản. Package PipelineALPR thực thi toàn bộ luồng xử lý AI, gồm năm sub-package hoạt động theo chuỗi: FrameProcessing giải mã và tiền xử lý khung hình từ video đầu vào; VehicleDetection phát hiện phương tiện trong từng khung hình; VehicleTracking theo dõi từng phương tiện xuyên suốt các khung hình liên tiếp; PlateDetection phát hiện và trích xuất vùng biển số; PlateOCR nhận dạng ký tự biển số; ResultAggregation tổng hợp kết quả từ nhiều khung hình để đưa ra kết quả cuối cùng cho từng phương tiện.
\paragraph{Tầng dữ liệu (Data Tier)}
Gồm hai package. Package Database lưu trữ toàn bộ dữ liệu có cấu trúc, gồm UserStore (thông tin tài khoản và phiên xác thực), RecognitionStore (kết quả các phiên nhận dạng và bản ghi từng phương tiện) và MonitorEventStore (thông tin các sự kiện được đánh dấu trong quá trình giám sát). Package ObjectStorage lưu trữ dữ liệu nhị phân, cụ thể là EvidenceImageStore chứa ảnh crop biển số và ảnh minh chứng phương tiện; package Database chỉ lưu URL tham chiếu đến các tài nguyên này thay vì nhúng trực tiếp ảnh vào document.
\section{Thiết kế chi tiết}
\subsection{Thiết kế giao diện}
Phần này có độ dài từ hai đến ba trang. Sinh viên đặc tả thông tin về màn hình mà ứng dụng của mình hướng tới, bao gồm độ phân giải màn hình, kích thước màn hình, số lượng màu sắc hỗ trợ, v.v. Tiếp đến, sinh viên đưa ra các thống nhất/chuẩn hóa của mình khi thiết kế giao diện như thiết kế nút, điều khiển, vị trí hiển thị thông điệp phản hồi, phối màu, v.v. Sau cùng sinh viên đưa ra một số hình ảnh minh họa thiết kế giao diện cho các chức năng quan trọng nhất. Lưu ý, sinh viên không nhầm lẫn giao diện thiết kế với giao diện của sản phẩm sau cùng.
\subsection{Thiết kế lớp}
Phần này có độ dài từ ba đến bốn trang. Sinh viên trình bày thiết kế chi tiết các thuộc tính và phương thức cho một số lớp chủ đạo/quan trọng nhất của ứng dụng (từ 2-4 lớp). Thiết kế chi tiết cho các lớp khác, nếu muốn trình bày, sinh viên đưa vào phần phụ lục.

Để minh họa thiết kế lớp, sinh viên thiết kế luồng truyền thông điệp giữa các đối tượng tham gia cho 2 đến 3 use case quan trọng nào đó bằng biểu đồ trình tự (hoặc biểu đồ giao tiếp).
\subsection{Thiết kế cơ sở dữ liệu}

Hệ thống sử dụng MongoDB làm cơ sở dữ liệu chính và Supabase Storage/S3-compatible storage làm kho lưu trữ ảnh evidence. MongoDB chỉ lưu dữ liệu có cấu trúc: tài khoản người dùng, phiên đăng nhập, phiên xử lý video, kết quả nhận dạng và các event được đánh dấu trong chế độ monitor. Ảnh crop phương tiện, ảnh crop biển số và các frame evidence được upload lên bucket \texttt{evidence}; trong MongoDB chỉ lưu URL trỏ tới các object này.

Thiết kế này bám theo cách dữ liệu được sinh ra trong pipeline. Một phiên xử lý video có thể sinh nhiều phương tiện được theo dõi, mỗi phương tiện có một biển số cuối cùng và nhiều crop biển số theo thời gian. Nếu lưu ảnh nhị phân trực tiếp trong document, kích thước document tăng nhanh và khó mở rộng khi track kéo dài. Vì vậy, hệ thống tách metadata và ảnh evidence: MongoDB giữ khóa truy vấn, kết quả OCR và các URL; object storage giữ dữ liệu ảnh.

Về mô hình NoSQL, hệ thống kết hợp hai cách lưu. Các thực thể có vòng đời riêng và cần truy vấn độc lập được lưu thành collection riêng, ví dụ \texttt{users}, \texttt{auth\_sessions}, \texttt{recognition\_sessions}, \texttt{recognition\_records} và \texttt{events}. Các dữ liệu nhỏ, chỉ có ý nghĩa trong ngữ cảnh document cha, được nhúng trực tiếp: \texttt{PlateFrame} nằm trong \texttt{RecognitionRecord}, còn \texttt{EventVehicle} nằm trong \texttt{MonitorEvent}. Cách này giảm số lần truy vấn khi giao diện mở chi tiết một kết quả nhận dạng hoặc một event monitor.

\subsubsection{Biểu đồ thực thể liên kết logic}

Mặc dù MongoDB không ràng buộc khóa ngoại như cơ sở dữ liệu quan hệ, biểu đồ E-R vẫn được dùng để mô tả quan hệ logic giữa các document chính. Các mũi tên nét liền biểu diễn quan hệ tham chiếu bằng khóa như \texttt{user\_id}, \texttt{session\_id}; mũi tên nét đứt biểu diễn quan hệ nhúng document con hoặc quan hệ logic không bắt buộc có khóa ngoại vật lý.

\begin{figure}[H]
    \centering
    \resizebox{\textwidth}{!}{%
    \begin{tikzpicture}[
        entity/.style={
            draw,
            rounded corners,
            align=left,
            text width=3.35cm,
            inner sep=5pt,
            font=\scriptsize
        },
        embedded/.style={
            entity,
            dashed
        },
        external/.style={
            entity,
            draw=gray
        },
        rel/.style={-{Latex[length=2mm]}, thick},
        embedrel/.style={-{Latex[length=2mm]}, thick, dashed},
        every edge quotes/.style={font=\scriptsize, fill=white, inner sep=1pt}
    ]
        \node[entity] (user) at (0, 4.8) {
            \textbf{User}\\
            \textit{collection: \texttt{users}}\\
            \texttt{\_id}, \texttt{email}\\
            \texttt{name}, \texttt{password\_hash}\\
            \texttt{role}, \texttt{is\_active}
        };

        \node[entity] (auth) at (4.6, 4.8) {
            \textbf{Auth}\\
            \textbf{Session}\\
            \textit{collection: \texttt{auth\_sessions}}\\
            \texttt{session\_id}\\
            \texttt{user\_id}\\
            \texttt{expires\_at}, \texttt{revoked}
        };

        \node[entity] (session) at (0, 2.1) {
            \textbf{Recognition}\\
            \textbf{Session}\\
            \textit{collection:}\\
            \texttt{recognition\_sessions}\\
            \texttt{session\_id}, \texttt{user\_id}\\
            \texttt{source\_filename}\\
            \texttt{status}, \texttt{ocr\_backend}
        };

        \node[entity] (record) at (4.6, 2.1) {
            \textbf{Recognition}\\
            \textbf{Record}\\
            \textit{collection:}\\
            \texttt{recognition\_records}\\
            \texttt{session\_id}, \texttt{track\_id}\\
            \texttt{plate\_text}\\
            \texttt{plate\_text\_confidence}
        };

        \node[embedded] (frame) at (9.2, 2.1) {
            \textbf{Plate}\\
            \textbf{Frame}\\
            \textit{embedded document}\\
            \texttt{frame\_index}\\
            \texttt{quality\_score}\\
            \texttt{image\_url}\\
            \texttt{ocr\_text}
        };

        \node[entity] (event) at (0, -1.0) {
            \textbf{Monitor}\\
            \textbf{Event}\\
            \textit{collection: \texttt{events}}\\
            \texttt{event\_id}\\
            \texttt{session\_id}\\
            \texttt{source\_type}, \texttt{source\_ref}\\
            \texttt{marked\_at}, \texttt{status}
        };

        \node[embedded] (eventvehicle) at (4.6, -1.0) {
            \textbf{Event}\\
            \textbf{Vehicle}\\
            \textit{embedded document}\\
            \texttt{track\_id}\\
            \texttt{plate\_text}\\
            \texttt{vehicle\_class}\\
            \texttt{plate\_image\_url}
        };

        \node[external] (evidence) at (9.2, -1.0) {
            \textbf{Evidence}\\
            \textbf{Object}\\
            \textit{Supabase/S3}\\
            \texttt{bucket=evidence}\\
            \texttt{object\_path}\\
            \texttt{public\_url}
        };

        \draw[rel] (user) -- node[above] {1-n} (auth);
        \draw[rel] (user) -- node[left] {1-n} (session);
        \draw[rel] (session) -- node[above] {1-n} (record);
        \draw[embedrel] (record) -- node[above] {embeds} (frame);
        \draw[embedrel] (event) -- node[above] {embeds} (eventvehicle);
        \draw[embedrel] (session) -- node[left] {\texttt{session\_id}} (event);
        \draw[embedrel] (record) |- node[pos=0.28, right] {URL} (evidence);
        \draw[embedrel] (frame) |- node[pos=0.25, right] {URL} (evidence);
        \draw[embedrel] (eventvehicle) -- node[above] {URL} (evidence);
    \end{tikzpicture}%
    }
    \caption{Biểu đồ E-R logic và ánh xạ NoSQL của hệ thống}
    \label{fig:database-er}
\end{figure}

Các thực thể chính trong hệ thống gồm:

\begin{itemize}
    \item \textbf{User}: tài khoản người dùng của hệ thống. Collection triển khai là \texttt{users}.
    \item \textbf{AuthSession}: phiên đăng nhập phía server, dùng cùng JWT/HttpOnly cookie để xác thực request. Collection triển khai là \texttt{auth\_sessions}.
    \item \textbf{RecognitionSession}: một phiên xử lý video hoặc nguồn dữ liệu nhận dạng. Collection triển khai là \texttt{recognition\_sessions}.
    \item \textbf{RecognitionRecord}: kết quả cuối cùng của một phương tiện trong một phiên xử lý, gồm biển số, độ tin cậy OCR, lớp phương tiện, track id và evidence. Collection triển khai là \texttt{recognition\_records}.
    \item \textbf{PlateFrame}: một crop biển số trong quá trình theo dõi. Đây là document con, được nhúng trong \texttt{RecognitionRecord}.
    \item \textbf{MonitorEvent}: event được người dùng đánh dấu trong chế độ monitor. Trong project, collection hiện tại là \texttt{events}, không còn dùng \texttt{incidents}.
    \item \textbf{EventVehicle}: một phương tiện được phát hiện trong cửa sổ thời gian của event. Đây là document con, được nhúng trong \texttt{MonitorEvent}.
    \item \textbf{EvidenceObject}: ảnh evidence lưu trong Supabase Storage/S3. Đây không phải collection MongoDB; MongoDB chỉ lưu URL.
\end{itemize}

Một \texttt{User} có nhiều \texttt{AuthSession} và nhiều \texttt{RecognitionSession}; cả hai quan hệ này đều dùng trường \texttt{user\_id}. Một \texttt{RecognitionSession} có nhiều \texttt{RecognitionRecord}; mỗi record dùng \texttt{session\_id} để trỏ về phiên xử lý và dùng cặp \texttt{(session\_id, track\_id)} để tránh lưu trùng cùng một track. Đối với monitor, \texttt{MonitorEvent.session\_id} dùng để gom các event thuộc cùng một phiên monitor hoặc cùng nguồn upload/live; đây là quan hệ logic trong ứng dụng, không phải khóa ngoại bắt buộc của MongoDB.

\texttt{PlateFrame} không được tách thành collection riêng vì frame chỉ có ý nghĩa khi đi cùng một record nhận dạng. Mỗi \texttt{RecognitionRecord} chứa một \texttt{best\_plate\_frame} và một mảng \texttt{track\_buffer}. Tương tự, danh sách \texttt{EventVehicle} được nhúng trong \texttt{MonitorEvent} vì giao diện event luôn cần đọc cùng lúc trạng thái event, danh sách xe và URL ảnh evidence của từng xe.

\subsubsection{Ánh xạ sang mô hình MongoDB}

Từ mô hình logic trên, hệ thống triển khai thành các collection/document như Bảng \ref{tab:mongodb-collections}. Các tên trong bảng được căn theo tên model và collection thật trong source code.

\begin{table}[H]
\centering
\caption{Ánh xạ thực thể sang collection/document MongoDB}
\label{tab:mongodb-collections}
\renewcommand{\arraystretch}{1.12}
\begin{tabular}{|>{\raggedright\arraybackslash}p{3.0cm}|>{\raggedright\arraybackslash}p{3.6cm}|>{\raggedright\arraybackslash}p{7.6cm}|}
\hline
\textbf{Thực thể logic} & \textbf{Collection/Document} & \textbf{Mô tả thiết kế} \\
\hline
User & \texttt{users} & Collection độc lập. Lưu thông tin tài khoản, email, tên hiển thị, mật khẩu đã băm, vai trò và trạng thái hoạt động. \\
\hline
AuthSession & \texttt{auth\_sessions} & Collection độc lập. Tham chiếu đến \texttt{users} bằng \texttt{user\_id}; dùng để kiểm tra phiên đăng nhập và thu hồi session. \\
\hline
RecognitionSession & \shortstack[l]{\texttt{recognition\_}\\\texttt{sessions}} & Collection độc lập. Lưu metadata của phiên xử lý như file nguồn, trạng thái, tổng số record, số frame đã xử lý, chế độ tiền xử lý và OCR backend. \\
\hline
RecognitionRecord & \shortstack[l]{\texttt{recognition\_}\\\texttt{records}} & Collection độc lập. Tham chiếu logic đến phiên xử lý bằng \texttt{session\_id}; lưu kết quả biển số cuối cùng của từng track. \\
\hline
PlateFrame & Embedded document & Được nhúng trong \texttt{RecognitionRecord} dưới dạng \texttt{best\_plate\_frame} và mảng \texttt{track\_buffer}. \\
\hline
MonitorEvent & \texttt{events} & Collection độc lập. Lưu event được đánh dấu trong chế độ monitor, gồm nguồn, khoảng thời gian, trạng thái xử lý và danh sách phương tiện trong event. \\
\hline
EventVehicle & Embedded document & Được nhúng trong \texttt{MonitorEvent} dưới dạng mảng \texttt{vehicles}. \\
\hline
EvidenceObject & Supabase/S3 object & Không phải collection MongoDB. File ảnh được lưu trong bucket object storage; MongoDB chỉ lưu URL trong các trường \texttt{image\_url}, \texttt{vehicle\_thumbnail\_url}, \texttt{plate\_image\_url}, \texttt{vehicle\_image\_url}. \\
\hline
\end{tabular}
\end{table}

\subsubsection{Thiết kế các collection chính}

\textbf{Collection \texttt{users}.} Collection này lưu tài khoản người dùng. Các trường chính gồm \texttt{\_id}, \texttt{email}, \texttt{name}, \texttt{password\_hash}, \texttt{role}, \texttt{is\_active}, \texttt{created\_at} và \texttt{updated\_at}. Trường \texttt{email} được đặt unique index để không có hai tài khoản dùng cùng một email. Hệ thống không lưu mật khẩu gốc.

\textbf{Collection \texttt{auth\_sessions}.} Collection này lưu phiên đăng nhập phía server. Mỗi document gồm \texttt{session\_id}, \texttt{user\_id}, \texttt{expires\_at}, \texttt{revoked}, \texttt{created\_at} và \texttt{updated\_at}. Khi đăng nhập, backend tạo một session, sinh JWT chứa \texttt{user\_id} và \texttt{session\_id}, sau đó gửi về trình duyệt bằng HttpOnly cookie. Khi xác thực request, backend giải mã JWT và đối chiếu với document trong \texttt{auth\_sessions}.

\textbf{Collection \texttt{recognition\_sessions}.} Collection này biểu diễn một phiên xử lý video hoặc nguồn nhận dạng. Các trường chính gồm \texttt{session\_id}, \texttt{user\_id}, \texttt{source\_filename}, \texttt{source\_type}, \texttt{status}, \texttt{total\_records}, \texttt{processed\_frames}, \texttt{preprocess\_mode}, \texttt{ocr\_backend}, \texttt{error\_message}, \texttt{created\_at} và \texttt{updated\_at}. Khi người dùng upload video, backend tạo session ở trạng thái \texttt{processing}; khi pipeline kết thúc, session được cập nhật thành \texttt{completed} hoặc \texttt{failed}.

\textbf{Collection \texttt{recognition\_records}.} Collection này lưu kết quả cuối cùng cho từng phương tiện được theo dõi. Các trường định danh gồm \texttt{session\_id}, \texttt{user\_id}, \texttt{track\_id}, \texttt{vehicle\_track\_id}, \texttt{plate\_track\_id} và \texttt{vehicle\_class}. Các trường OCR gồm \texttt{plate\_text}, \texttt{plate\_text\_confidence}, \texttt{ocr\_vote\_summary} và \texttt{ocr\_method}. Các trường thời gian gồm \texttt{first\_seen\_frame}, \texttt{last\_seen\_frame}, \texttt{created\_at} và \texttt{updated\_at}.

Trong \texttt{recognition\_records}, phần evidence được tổ chức thành ba nhóm. Trường \texttt{vehicle\_thumbnail\_url} lưu URL ảnh phương tiện. Trường \texttt{best\_plate\_frame} lưu crop biển số tốt nhất dưới dạng \texttt{PlateFrame}. Trường \texttt{track\_buffer} là mảng \texttt{PlateFrame}, mỗi phần tử gồm \texttt{frame\_index}, \texttt{quality\_score}, \texttt{image\_url}, \texttt{timestamp\_ms}, \texttt{plate\_bbox}, \texttt{detection\_confidence}, \texttt{ocr\_text} và \texttt{ocr\_confidence}. Đây là dữ liệu phục vụ màn hình xem evidence và kiểm tra lại các frame đã đóng góp vào kết quả OCR.

\textbf{Collection \texttt{events}.} Collection này thay thế mô hình \texttt{incidents} cũ. Mỗi document \texttt{MonitorEvent} lưu một cửa sổ thời gian được người dùng đánh dấu trong chế độ monitor. Các trường chính gồm \texttt{event\_id}, \texttt{session\_id}, \texttt{source\_type}, \texttt{source\_ref}, \texttt{marked\_at}, \texttt{window\_start\_sec}, \texttt{window\_end\_sec}, \texttt{duration\_sec}, \texttt{status}, \texttt{vehicles}, \texttt{total\_vehicles}, \texttt{processing\_ms}, \texttt{error\_message}, \texttt{created\_at} và \texttt{updated\_at}. Trường \texttt{vehicles} là mảng \texttt{EventVehicle}. Mỗi phần tử lưu \texttt{track\_id}, \texttt{vehicle\_track\_id}, \texttt{plate\_track\_id}, \texttt{plate\_text}, \texttt{plate\_text\_confidence}, \texttt{chars}, \texttt{vehicle\_class}, \texttt{plate\_image\_url}, \texttt{vehicle\_image\_url}, \texttt{ocr\_method}, \texttt{ocr\_frames}, \texttt{first\_seen\_frame} và \texttt{last\_seen\_frame}.

\subsubsection{Thiết kế lưu trữ evidence bằng Supabase Storage/S3}

Các ảnh evidence được lưu vào bucket \texttt{evidence}. Trong luồng xử lý video thông thường, hệ thống lưu ảnh theo mẫu đường dẫn:

\begin{itemize}
    \item \texttt{\{session\_id\}/track\_\{track\_id\}\_frame\_\{index\}.jpg}: các frame trong track buffer;
    \item \texttt{\{session\_id\}/plate\_\{track\_id\}.jpg}: ảnh biển số tốt nhất;
    \item \texttt{\{session\_id\}/vehicle\_\{track\_id\}.jpg}: ảnh crop phương tiện.
\end{itemize}

Trong luồng monitor event, ảnh evidence được lưu theo mẫu:

\begin{itemize}
    \item \texttt{events/\{event\_id\}/plate\_\{id\}.jpg};
    \item \texttt{events/\{event\_id\}/vehicle\_\{id\}.jpg}.
\end{itemize}

Sau khi upload ảnh, backend lấy public URL và lưu URL đó vào MongoDB. Nếu upload không thành công, URL có thể rỗng nhưng kết quả nhận dạng vẫn được lưu, giúp pipeline không phụ thuộc tuyệt đối vào object storage. Nếu triển khai trong môi trường cần bảo mật evidence chặt hơn, bucket có thể chuyển sang private và backend sinh signed URL có thời hạn thay vì lưu public URL.

\subsubsection{Thiết kế chỉ mục}

Các chỉ mục được tạo trong hàm khởi tạo MongoDB để phục vụ đúng các truy vấn chính của ứng dụng: đăng nhập bằng email, kiểm tra session, liệt kê lịch sử theo người dùng, lấy record theo phiên và lấy event gần nhất. Bảng \ref{tab:mongodb-indexes} liệt kê các chỉ mục chính.

\begin{table}[H]
\centering
\caption{Các chỉ mục MongoDB chính}
\label{tab:mongodb-indexes}
\renewcommand{\arraystretch}{1.12}
\begin{tabular}{|>{\raggedright\arraybackslash}p{3.4cm}|>{\raggedright\arraybackslash}p{5.6cm}|>{\raggedright\arraybackslash}p{5.4cm}|}
\hline
\textbf{Collection} & \textbf{Chỉ mục} & \textbf{Mục đích} \\
\hline
\texttt{users} & \texttt{email} unique; \newline \texttt{created\_at} desc & Đảm bảo email không trùng, tăng tốc đăng nhập và hỗ trợ liệt kê tài khoản mới. \\
\hline
\texttt{auth\_sessions} & \texttt{session\_id} unique; \newline \texttt{user\_id}; \newline \texttt{expires\_at} & Kiểm tra phiên đăng nhập, truy vết phiên theo người dùng và hỗ trợ xử lý phiên hết hạn. \\
\hline
\shortstack[l]{\texttt{recognition\_}\\\texttt{sessions}} & \texttt{session\_id} unique; \newline \texttt{(user\_id, created\_at)}; \newline \texttt{status}; \newline \texttt{created\_at} desc & Lấy phiên theo mã, liệt kê lịch sử gần nhất của người dùng và lọc theo trạng thái xử lý. \\
\hline
\shortstack[l]{\texttt{recognition\_}\\\texttt{records}} & \texttt{session\_id}; \newline \texttt{(user\_id, session\_id)}; \newline \texttt{plate\_text}; \newline \texttt{vehicle\_track\_id}; \newline \texttt{plate\_track\_id}; \newline \texttt{vehicle\_class}; \newline \texttt{created\_at} desc; \newline \texttt{(session\_id, track\_id)} unique & Lấy record của một phiên, lấy record theo track, tìm kiếm biển số, lọc theo mã track phương tiện/biển số và tránh lưu trùng cùng một track trong một session. \\
\hline
\texttt{events} & \texttt{event\_id} unique; \newline \texttt{session\_id}; \newline \texttt{marked\_at} desc; \newline \texttt{status}; \newline \texttt{source\_type} & Lấy chi tiết event, liệt kê event theo phiên monitor, theo nguồn và theo thời gian đánh dấu. \\
\hline
\end{tabular}
\end{table}

\subsubsection{Luồng truy vấn chính}

Thiết kế collection và chỉ mục được chọn dựa trên các luồng truy vấn chính của ứng dụng. Khi người dùng đăng nhập, hệ thống tìm tài khoản theo \texttt{email}, kiểm tra mật khẩu và tạo document trong \texttt{auth\_sessions}. Khi người dùng mở lịch sử, frontend gọi API lấy danh sách \texttt{recognition\_sessions} theo \texttt{user\_id}, sắp xếp giảm dần theo \texttt{created\_at}. Khi người dùng chọn một phiên, backend lấy các \texttt{recognition\_records} theo \texttt{session\_id} và \texttt{user\_id}. Khi người dùng mở track buffer, backend lấy một record theo cặp \texttt{(session\_id, track\_id)}.

Đối với monitor, khi người dùng đánh dấu một khoảng thời gian, backend tạo \texttt{event\_id}, chạy pipeline ALPR trên cửa sổ đó, phát tiến độ qua SSE tại \texttt{/monitor/\{session\_id\}/events/stream}, rồi lưu kết quả vào collection \texttt{events}. Khi cần xem lại, backend cung cấp API \texttt{/events/\{event\_id\}} để lấy một event và \texttt{/events} để liệt kê event theo \texttt{session\_id}, \texttt{source\_type} hoặc giới hạn số lượng. Vì danh sách \texttt{EventVehicle} được nhúng trực tiếp trong document \texttt{MonitorEvent}, giao diện có thể hiển thị chi tiết event bằng một lần đọc document chính.

Tóm lại, thiết kế cơ sở dữ liệu hiện tại không còn sử dụng collection \texttt{incidents}. Toàn bộ luồng monitor được chuẩn hóa theo khái niệm \texttt{events}. MongoDB lưu metadata và kết quả nhận dạng, còn ảnh evidence được tách sang object storage để giảm kích thước document và giữ cho các truy vấn lịch sử, chi tiết record và chi tiết event đơn giản.


\section{Xây dựng ứng dụng}

\subsection*{Thư viện và công cụ sử dụng}

Hệ thống được phát triển bằng các công cụ, ngôn ngữ lập trình và thư viện chính trong Bảng \ref{table:tools_libraries}.

\renewcommand{\arraystretch}{1.15}
\begin{longtable}{p{3.2cm} p{4.2cm} p{2.6cm} p{4.8cm}}
\caption{Danh sách thư viện và công cụ sử dụng}
\label{table:tools_libraries} \\

\hline
\textbf{Mục đích} & \textbf{Công cụ/Thư viện} & \textbf{Phiên bản} & \textbf{Địa chỉ URL} \\ \hline
\endfirsthead

\hline
\textbf{Mục đích} & \textbf{Công cụ/Thư viện} & \textbf{Phiên bản} & \textbf{Địa chỉ URL} \\ \hline
\endhead

IDE lập trình & Visual Studio Code & 1.104.0 & \url{https://code.visualstudio.com/} \\ \hline
Quản lý mã nguồn & Git & 2.43.0 & \url{https://git-scm.com/} \\ \hline
Đóng gói ứng dụng & Docker & 28.5.1 & \url{https://www.docker.com/} \\ \hline
Ngôn ngữ backend & Python & 3.10 & \url{https://www.python.org/} \\ \hline
Runtime frontend & Node.js & 20 & \url{https://nodejs.org/} \\ \hline
Package manager & npm & 11.6.3 & \url{https://www.npmjs.com/} \\ \hline
Backend API & FastAPI & $\geq$ 0.111.0 & \url{https://fastapi.tiangolo.com/} \\ \hline
ASGI server & Uvicorn & $\geq$ 0.29.0 & \url{https://www.uvicorn.org/} \\ \hline
Học sâu & PyTorch & 2.4.1 & \url{https://pytorch.org/} \\ \hline
Xử lý ảnh & OpenCV Python & 4.10.0.84 & \url{https://opencv.org/} \\ \hline
Phát hiện đối tượng & Ultralytics YOLO & $\geq$ 8.0.0 & \url{https://docs.ultralytics.com/} \\ \hline
Theo dõi đối tượng & BoxMOT & 19.0.0 & \url{https://github.com/mikel-brostrom/boxmot} \\ \hline
Suy luận ONNX & ONNX Runtime & 1.20.0 & \url{https://onnxruntime.ai/} \\ \hline
Tính toán ma trận & NumPy & 1.26.4 & \url{https://numpy.org/} \\ \hline
Tính toán khoa học & SciPy & $\geq$ 1.11.0 & \url{https://scipy.org/} \\ \hline
Huấn luyện mô hình & Lightning & 2.x & \url{https://lightning.ai/} \\ \hline
Xử lý ảnh phụ trợ & Pillow & $\geq$ 10.0.0 & \url{https://python-pillow.org/} \\ \hline
Đọc cấu hình & PyYAML & $\geq$ 6.0.1 & \url{https://pyyaml.org/} \\ \hline
Xác thực mật khẩu & bcrypt & $\geq$ 4.1.0 & \url{https://github.com/pyca/bcrypt} \\ \hline
JWT & PyJWT & $\geq$ 2.8.0 & \url{https://pyjwt.readthedocs.io/} \\ \hline
Cơ sở dữ liệu & MongoDB & 7 & \url{https://www.mongodb.com/} \\ \hline
Driver MongoDB & Motor & $\geq$ 3.3.0 & \url{https://motor.readthedocs.io/} \\ \hline
Lưu trữ ảnh & Supabase & $\geq$ 2.5.0 & \url{https://supabase.com/} \\ \hline
Frontend framework & React & 18.3.1 & \url{https://react.dev/} \\ \hline
Điều hướng frontend & React Router DOM & 7.17.0 & \url{https://reactrouter.com/} \\ \hline
Build frontend & Vite & 5.4.21 & \url{https://vite.dev/} \\ \hline
Thiết kế giao diện & Tailwind CSS & 3.4.19 & \url{https://tailwindcss.com/} \\ \hline
Streaming video & MediaMTX & v1.19.1 & \url{https://github.com/bluenviron/mediamtx} \\ \hline
Xử lý video & FFmpeg & 6.1.1 & \url{https://ffmpeg.org/} \\ \hline
Kiểm thử & pytest & -- & \url{https://docs.pytest.org/} \\ \hline

\end{longtable}

\subsection{Frontend}

Frontend của hệ thống được xây dựng dưới dạng ứng dụng web một trang bằng React và Vite, mã nguồn đặt trong thư mục \texttt{web/src}. Giao diện không thực hiện suy luận mô hình trực tiếp, mà đóng vai trò tiếp nhận đầu vào từ người dùng, gửi yêu cầu tới backend, nhận tiến độ xử lý qua cơ chế streaming và hiển thị kết quả nhận dạng theo từng phương tiện hoặc từng event.

Ứng dụng được tổ chức quanh hai chế độ chính: chế độ xử lý video tải lên và chế độ giám sát. Tại \texttt{App.jsx}, trạng thái \texttt{mode} nhận hai giá trị \texttt{process} và \texttt{monitor}; người dùng chuyển chế độ bằng thành phần \texttt{SegmentedControl} trên thanh điều hướng. Cách tổ chức này giúp tách rõ hai luồng sử dụng: luồng \texttt{process} dùng cho đánh giá một video hoàn chỉnh, còn luồng \texttt{monitor} dùng cho việc kết nối camera hoặc mở video và đánh dấu một khoảng thời gian cần phân tích.

\begin{table}[H]
\centering
\renewcommand{\arraystretch}{1.15}
\begin{tabular}{|p{3.6cm}|p{9.8cm}|}
\hline
\textbf{Thành phần} & \textbf{Vai trò trong giao diện} \\ \hline
\texttt{AuthProvider}, \texttt{apiClient.js} & Quản lý trạng thái đăng nhập, gọi các API \texttt{/auth/me}, \texttt{/auth/login}, \texttt{/auth/register}, \texttt{/auth/logout} và tự động gắn CSRF token cho các yêu cầu ghi dữ liệu đi qua \texttt{apiFetch}. \\ \hline
\texttt{DashboardPage} & Điều phối trạng thái chính của màn hình làm việc: file video, \texttt{job\_id}, tiến độ xử lý, danh sách phương tiện hợp lệ, danh sách kết quả bị loại và lỗi trả về từ backend. \\ \hline
\texttt{SourcePanel} & Nhận file video, chọn chế độ tiền xử lý ảnh và backend OCR trước khi gửi yêu cầu \texttt{POST /upload}. \\ \hline
\texttt{MediaStage} & Hiển thị video gốc hoặc khung hình preview được backend gửi về trong quá trình xử lý. \\ \hline
\texttt{ResultsPanel}, \texttt{VehicleCard} & Hiển thị kết quả theo từng track phương tiện, gồm biển số tổng hợp, độ tin cậy, ảnh evidence và các khung hình trong track buffer. \\ \hline
\texttt{HistoryModal} & Truy vấn lại danh sách phiên và các record đã lưu thông qua \texttt{/sessions} và \texttt{/sessions/\{session\_id\}/records}. \\ \hline
\texttt{MonitorPage} & Quản lý phiên giám sát RTSP hoặc video upload, gửi lệnh đánh dấu event và nhận tiến độ phân tích event qua SSE. \\ \hline
\texttt{LiveViewer}, \texttt{UploadViewer}, \texttt{EventsPanel} & Hiển thị luồng camera/video, chọn khoảng thời gian cần phân tích và trình bày danh sách event đã đánh dấu. \\ \hline
\end{tabular}
\caption{Các thành phần frontend chính}
\label{table:frontend_components}
\end{table}

Trong luồng xử lý video tải lên, người dùng chọn file tại \texttt{SourcePanel}. Hook \texttt{useUpload} đóng gói file vào \texttt{FormData}, kèm theo \texttt{preprocess\_mode} và \texttt{ocr\_backend}, sau đó gọi \texttt{POST /upload}. Backend trả về \texttt{job\_id}; từ thời điểm này, hook \texttt{useStream} mở kết nối \texttt{EventSource} tới \texttt{/stream/\{job\_id\}} để nhận các sự kiện \texttt{progress}, \texttt{frame}, \texttt{vehicle}, \texttt{rejected\_vehicle}, \texttt{complete} và \texttt{error}. Dữ liệu nhận được được cập nhật vào state của \texttt{DashboardPage}, nhờ vậy giao diện có thể hiển thị tiến độ, khung hình đang xử lý và kết quả nhận dạng mà không cần tải lại trang.

Trong luồng giám sát, \texttt{MonitorPage} cho phép người dùng chọn một trong hai nguồn: RTSP camera hoặc video upload. Với RTSP, frontend gửi \texttt{rtsp\_url} tới \texttt{POST /monitor/live/connect} và nhận lại \texttt{session\_id}, \texttt{whep\_url}, \texttt{mjpeg\_url}. Với video upload, frontend gửi file tới \texttt{POST /monitor/upload} và nhận URL phát lại video. Khi người dùng đánh dấu một thời điểm hoặc một khoảng thời gian, frontend gọi \texttt{POST /monitor/\{session\_id\}/mark}; kết quả phân tích được đẩy ngược về qua \texttt{/monitor/\{session\_id\}/events/stream}. Cách này giúp thao tác đánh dấu event nhanh, trong khi quá trình ALPR vẫn chạy bất đồng bộ ở backend.

Về mặt trải nghiệm, frontend ưu tiên hiển thị theo tiến trình thực tế của pipeline. Các trạng thái rỗng, đang upload, đang xử lý, hoàn thành và lỗi được tách riêng bằng các thành phần giao diện như \texttt{Badge}, \texttt{Toast}, \texttt{Progress} và \texttt{EmptyState}. Điều này giúp người dùng theo dõi được hệ thống đang ở bước nào: chờ chọn nguồn, đang tải video, đang phân tích, đã có kết quả hoặc cần xử lý lỗi đầu vào.

\subsection{Backend và API}

Backend được xây dựng bằng FastAPI, mã nguồn chính đặt trong thư mục \texttt{api}. File \texttt{api/main.py} chỉ giữ vai trò định nghĩa route, kiểm tra đầu vào, quản lý hàng đợi streaming và điều phối job xử lý. Các phần tính toán nặng được tách sang thư mục \texttt{api/core}, gồm nạp mô hình, phát hiện phương tiện, phát hiện biển số, theo dõi đối tượng, OCR, chấm điểm chất lượng khung hình và tổng hợp kết quả theo track.

Khi ứng dụng khởi động, hàm \texttt{lifespan} của FastAPI nạp các mô hình một lần thông qua \texttt{load\_models()}, khởi tạo kết nối MongoDB nếu có cấu hình \texttt{MONGODB\_URI}, đồng thời bật tác vụ dọn dẹp các phiên monitor tạm thời. Cách nạp mô hình ở giai đoạn khởi động giúp tránh việc mỗi request upload lại phải khởi tạo YOLO, OCR hoặc ReID từ đầu.

\begin{table}[H]
\centering
\renewcommand{\arraystretch}{1.15}
\begin{tabular}{|p{3.6cm}|p{9.8cm}|}
\hline
\textbf{Nhóm API} & \textbf{Endpoint và chức năng} \\ \hline
Xác thực & \texttt{/auth/register}, \texttt{/auth/login}, \texttt{/auth/logout}, \texttt{/auth/me}, \texttt{/auth/csrf}. Nhóm API này tạo tài khoản, đăng nhập, đăng xuất, kiểm tra phiên người dùng và cấp CSRF token cho các yêu cầu ghi dữ liệu. \\ \hline
Xử lý video tải lên & \texttt{POST /upload} tạo job xử lý video; \newline \texttt{GET /stream/\{job\_id\}} trả tiến độ và kết quả qua SSE; \newline \texttt{GET /stream/\{job\_id\}/mjpeg} trả khung hình preview dạng MJPEG; \newline \texttt{GET /records/\{job\_id\}/\{track\_id\}} lấy lại kết quả chi tiết của một track. \\ \hline
Lịch sử nhận dạng & \texttt{GET /sessions}, \texttt{GET /sessions/\{session\_id\}}, \texttt{GET /sessions/\{session\_id\}/records}. Nhóm API này phục vụ màn hình lịch sử và chỉ trả dữ liệu thuộc về người dùng hiện tại. \\ \hline
Giám sát và event & \texttt{POST /monitor/upload}, \texttt{POST /monitor/live/connect}, \newline \texttt{POST /monitor/\{session\_id\}/mark}, \newline \texttt{GET /monitor/\{session\_id\}/events/stream}, \newline \texttt{GET /events/\{event\_id\}}, \texttt{GET /events}. Nhóm API này quản lý phiên monitor, đánh dấu event và truy vấn lại kết quả event đã lưu. \\ \hline
\end{tabular}
\caption{Các nhóm API chính của backend}
\label{table:backend_api_groups}
\end{table}

Luồng \texttt{POST /upload} được thiết kế theo cơ chế bất đồng bộ. Backend kiểm tra định dạng file theo phần mở rộng \texttt{.mp4}, \texttt{.avi}, \texttt{.webm}, \texttt{.mov}, \texttt{.mkv}, kiểm tra kiểu nội dung, dung lượng tối đa và file rỗng. Sau đó hệ thống sinh \texttt{job\_id}, tạo hàng đợi SSE, ghi video vào file tạm và đưa hàm \texttt{run\_job} vào executor. Mỗi job được gắn với \texttt{user\_id}; các endpoint \texttt{/stream/\{job\_id\}} và \texttt{/stream/\{job\_id\}/mjpeg} đều kiểm tra chủ sở hữu job trước khi trả dữ liệu.

Trong quá trình chạy pipeline, backend đẩy các bản tin JSON vào hàng đợi của job. Các bản tin này gồm tiến độ theo frame, khung hình preview, kết quả phương tiện hợp lệ, kết quả bị loại do không đạt điều kiện hậu xử lý, trạng thái hoàn thành và lỗi. FastAPI chuyển các bản tin này thành luồng \texttt{text/event-stream}, nhờ đó frontend có thể cập nhật giao diện gần thời gian thực mà không cần polling liên tục.

Nhóm API xác thực được triển khai trong \texttt{api/auth.py}. Mật khẩu người dùng được băm bằng bcrypt, phiên đăng nhập được lưu trong collection \texttt{auth\_sessions}, còn trình duyệt nhận JWT trong cookie \texttt{HttpOnly}. Với các API ghi dữ liệu thuộc luồng xử lý video chính, backend yêu cầu CSRF token trong header \texttt{X-CSRF-Token}. Cách triển khai này tách dữ liệu xác thực khỏi mã frontend và tránh lưu token truy cập trực tiếp trong \texttt{localStorage}.

Đối với chức năng giám sát, file \texttt{api/routes\_monitor.py} quản lý hai loại phiên: \texttt{upload} và \texttt{live}. Phiên \texttt{upload} lưu video tạm để người dùng chọn khoảng thời gian cần phân tích; phiên \texttt{live} tạo \texttt{LiveSession}, kết nối nguồn RTSP và phát lại qua MediaMTX. Khi người dùng đánh dấu event, backend tạo \texttt{event\_id}, dựng \texttt{FrameSource} tương ứng với nguồn video, rồi gọi \texttt{run\_event} trong \texttt{api/core/event\_analyzer.py}. Kết quả event được phát về frontend qua SSE và được lưu vào collection \texttt{events}.

Lớp lưu trữ được tách khỏi route xử lý chính. MongoDB lưu thông tin người dùng, phiên xử lý, record nhận dạng và event monitor thông qua các hàm trong \texttt{api/database/mongodb.py}. Supabase Storage được sử dụng cho ảnh evidence của track và event để tránh lưu ảnh trực tiếp trong document MongoDB. Nhờ cách tách này, backend API giữ vai trò điều phối, còn dữ liệu lớn và metadata được lưu ở các thành phần phù hợp hơn.

\subsection{Kết quả đạt được}
Sinh viên trước tiên mô tả kết quả đạt được của mình là gì, ví dụ như các sản phẩm được đóng gói là gì, bao gồm những thành phần nào, ý nghĩa, vai trò?

Sinh viên cần thống kê các thông tin về ứng dụng của mình như: số dòng code, số lớp, số gói, dung lượng toàn bộ mã nguồn, dung lượng của từng sản phẩm đóng gói, v.v. Tương tự như phần liệt kê về công cụ sử dụng, sinh viên cũng nên dùng bảng để mô tả phần thông tin thống kê này.

\subsection{Minh họa các chức năng chính}
Sinh viên lựa chọn và đưa ra màn hình cho các chức năng chính, quan trọng, và thú vị nhất. Mỗi giao diện cần phải có lời giải thích ngắn gọn. Khi giải thích, sinh viên có thể kết hợp với các chú thích ở trong hình ảnh giao diện.

\section{Kiểm thử}
Phần này có độ dài từ hai đến ba trang. Sinh viên thiết kế các trường hợp kiểm thử cho hai đến ba chức năng quan trọng nhất. Sinh viên cần chỉ rõ các kỹ thuật kiểm thử đã sử dụng. Chi tiết các trường hợp kiểm thử khác, nếu muốn trình bày, sinh viên đưa vào phần phụ lục.
Sinh viên sau cùng tổng kết về số lượng các trường hợp kiểm thử và kết quả kiểm thử. Sinh viên cần phân tích lý do nếu kết quả kiểm thử không đạt.
\section{Triển khai}

Ứng dụng được triển khai theo mô hình tách riêng frontend, backend xử lý AI và các dịch vụ lưu trữ. Frontend được triển khai trên Vercel, backend được đóng gói bằng Docker và chạy trên RunPod GPU. Cơ sở dữ liệu và lưu trữ ảnh kết quả sử dụng các dịch vụ bên ngoài để giảm tải cho server xử lý chính.

\begin{table}[H]
\centering
\begin{tabular}{p{4cm} p{10cm}}
\hline
\textbf{Thành phần} & \textbf{Cấu hình triển khai} \\ \hline
Frontend & React/Vite, triển khai trên Vercel \\ \hline
Backend API & FastAPI, chạy trên cổng 8000 \\ \hline
Server xử lý AI & RunPod GPU, NVIDIA RTX 3090 \\ \hline
Môi trường GPU & PyTorch 2.4.1, CUDA 12.1, cuDNN 9 \\ \hline
Streaming video & MediaMTX, sử dụng WebRTC/RTSP \\ \hline
Cơ sở dữ liệu & MongoDB Atlas \\ \hline
Lưu trữ ảnh kết quả & Supabase Storage \\ \hline
Giới hạn dung lượng upload & 512 MB \\ \hline
\end{tabular}
\caption{Cấu hình triển khai ứng dụng}
\label{table:deployment_config}
\end{table}

Trong môi trường thử nghiệm cục bộ, hệ thống được chạy bằng Docker Compose gồm bốn dịch vụ chính: backend API, frontend web, MediaMTX và MongoDB. Backend mở cổng \texttt{8000}, frontend mở cổng \texttt{3000}, MongoDB mở cổng \texttt{27017}, MediaMTX sử dụng các cổng \texttt{8554}, \texttt{8889}, \texttt{8189/udp} và \texttt{9997}. Cấu hình này phục vụ kiểm thử chức năng đăng nhập, tải video, xử lý nhận dạng và xem kết quả trực tiếp.

Trong môi trường triển khai thử nghiệm trên server GPU, backend và MediaMTX được đóng gói trong cùng một Docker image. Backend đảm nhiệm các bước phát hiện phương tiện, phát hiện biển số, tracking, OCR và trả kết quả qua API. MediaMTX phục vụ luồng video trực tiếp, hỗ trợ hiển thị kết quả giám sát trên giao diện web. Các mô hình như detector phương tiện, detector biển số, OCR và ReID được đặt trên network volume của RunPod.
\end{document}
